\begin{ouverture}[calcul]{9}{Les décimaux jusqu’aux millièmes}{Partie 2 · Fractions et nombres décimaux}
{Millièmes · Comparaisons · Encadrements · Arrondis}
\objectif{Lire, décomposer et écrire des décimaux jusqu’aux millièmes.}
\objectif{Comparer, encadrer, intercaler des décimaux.}
\objectif{Donner la partie entière et l’arrondi à l’entier.}
\prerequis{Dixièmes et centièmes du CM1.}
\end{ouverture}
\pagestyle{otMain}\essentiel\fontsize{11}{15}\selectfont
\phantomsection\label{cours-9-1}
\essSection{1}{Le rang des millièmes}
\begin{definition}Partager un centième en 10 donne des millièmes : $\frac1{1\,000}$. Une unité vaut 10 dixièmes, 100 centièmes ou $1\,000$ millièmes. L’écriture $4{,}375$ code $4+\frac3{10}+\frac7{100}+\frac5{1\,000}$.\end{definition}\begin{center}\renewcommand{\arraystretch}{1.35}\begin{tabular}{c|cccc}\toprule
Nombre & unités & dixièmes & centièmes & millièmes\\\midrule
$4{,}375$ & 4 & 3 & 7 & 5\\
$6{,}04$ & 6 & 0 & 4 & \\
$0{,}208$ & 0 & 2 & 0 & 8\\
\bottomrule\end{tabular}\end{center}
\begin{exemple}$\frac{2\,375}{1\,000}=2+\frac{375}{1\,000}=2{,}375$. Et $0{,}208$ se lit « 208 millièmes » : $\frac{208}{1\,000}$, soit 2 dixièmes et 8 millièmes, sans centième.\end{exemple}
\begin{savoirfaire}[Des millièmes bien réels]Les millièmes mesurent finement : $1$ m $=1\,000$ mm, donc 1 mm $=0{,}001$ m ; une masse de 1 g vaut $0{,}001$ kg. Lire « 3 m et 45 mm » : $3{,}045$ m.\end{savoirfaire}
\clearpage
\phantomsection\label{cours-9-2}
\essSection{2}{Comparer et arrondir}
\begin{propriete}Comparer rang par rang, après la partie entière : $6{,}4>6{,}385$, car 4 dixièmes dépassent 3 dixièmes ; le nombre de chiffres ne décide pas. Compléter par des zéros aide : $6{,}400$ contre $6{,}385$.\end{propriete}
\begin{savoirfaire}[Encadrer et intercaler]Intercaler est toujours possible : entre $5{,}27$ et $5{,}28$, penser en millièmes : $5{,}275$ convient. Encadrer au dixième : $8{,}236$ est entre $8{,}2$ et $8{,}3$.\end{savoirfaire}
\begin{exemple}Arrondi à l’entier : regarder les dixièmes. $7{,}493\to7$ ; $7{,}51\to8$ ; $12{,}5\to13$ par convention à mi-chemin. Sur la demi-droite, l’arrondi est l’entier le plus proche.\end{exemple}
\begin{attention}$0{,}5$ ; $0{,}50$ ; $0{,}500$ désignent le même nombre : les zéros finaux après la virgule ne changent rien. En revanche $0{,}05$ et $0{,}005$ sont dix fois plus petits l’un que l’autre.\end{attention}

\clearpage\rubrique{Découverte}
\begin{blocLivre}{ex-9-1}\begin{exercice}[1]{Lire les rangs}
Dans $4{,}375$, donner le chiffre des centièmes et le chiffre des millièmes, puis décomposer le nombre.
\end{exercice}\end{blocLivre}
\begin{blocLivre}{ex-9-2}\begin{exercice}[1]{Des fractions aux virgules}
Écrire avec une virgule : $\frac{2\,375}{1\,000}$ ; $\frac{208}{1\,000}$ ; $\frac{45}{1\,000}$.
\end{exercice}\end{blocLivre}
\begin{blocLivre}{ex-9-3}\begin{exercice}[1]{Des virgules aux fractions}
Décomposer $6{,}04$ et $0{,}208$ avec des fractions décimales.
\end{exercice}\end{blocLivre}
\begin{blocLivre}{ex-9-4}\begin{exercice}[1]{Mètres et millimètres}
Écrire en mètres : 3 m et 45 mm ; 2 m et 7 mm ; 605 mm.
\end{exercice}\end{blocLivre}
\begin{blocLivre}{ex-9-5}\begin{exercice}[1]{Le piège des longueurs d’écriture}
Comparer $6{,}4$ et $6{,}385$, puis $0{,}72$ et $0{,}702$.
\end{exercice}\end{blocLivre}

\rubrique{Entraînement}
\begin{blocLivre}{ex-9-6}\begin{exercice}[2]{Ranger}
Ranger dans l’ordre croissant : $3{,}205$ ; $3{,}25$ ; $3{,}052$ ; $3{,}5$ ; $3{,}055$.
\end{exercice}\end{blocLivre}
\begin{blocLivre}{ex-9-7}\begin{exercice}[2]{Encadrer au dixième}
Encadrer $8{,}236$ entre deux dixièmes consécutifs, puis entre deux centièmes consécutifs.
\end{exercice}\end{blocLivre}
\begin{blocLivre}{ex-9-8}\begin{exercice}[2]{Intercaler}
Donner un nombre strictement compris entre $5{,}27$ et $5{,}28$, puis entre $0{,}999$ et 1.
\end{exercice}\end{blocLivre}
\begin{blocLivre}{ex-9-9}\begin{exercice}[2]{Arrondir}
Donner l’arrondi à l’entier de $7{,}493$ ; $7{,}51$ ; $12{,}5$.
\end{exercice}\end{blocLivre}
\begin{blocLivre}{ex-9-10}\begin{exercice}[2]{Sur la demi-droite}
La demi-droite est graduée de $0{,}01$ en $0{,}01$. Lire A et B, puis placer $C(4{,}235)$.
\label{fig-ex-9-10}\figureGeo{dec-10}
\end{exercice}\end{blocLivre}

\rubrique{Problèmes}
\begin{blocLivre}{ex-9-11}\begin{exercice}[3]{Les performances}
Au saut en longueur, Aya saute $3{,}205$ m, Lou $3{,}25$ m et Mia $3{,}209$ m. Donner le podium, puis l’écart entre les deux premières.
\end{exercice}\end{blocLivre}
\begin{blocLivre}{ex-9-12}\begin{exercice}[3]{Le même nombre ?}
Parmi ces écritures, regrouper celles qui désignent le même nombre : $0{,}5$ ; $0{,}050$ ; $0{,}500$ ; $\frac5{10}$ ; $\frac{50}{1\,000}$.
\end{exercice}\end{blocLivre}
\begin{blocLivre}{ex-9-13}\begin{exercice}[3]{La bonne précision}
Un robot découpe des pièces de $0{,}125$ m. Exprimer cette longueur en millimètres, puis dire combien de pièces entières il découpe dans une barre de 1 m.
\end{exercice}\end{blocLivre}
\begin{blocLivre}{ex-9-14}\begin{exercice}[3]{L’arrondi trompeur}
Deux nombres différents peuvent-ils avoir le même arrondi à l’entier ? Donner un exemple avec des millièmes, puis expliquer pourquoi l’arrondi ne suffit pas pour comparer.
\end{exercice}\end{blocLivre}
\begin{blocLivre}{ex-9-15}\begin{exercice}[3]{Zéro virgule neuf cent quatre-vingt-dix-neuf}
Nino affirme : « $0{,}999$ est plus grand que $0{,}99$ car il a plus de chiffres ; et il est plus grand que 1 car 999 est grand. » Démêler le vrai du faux.
\end{exercice}\end{blocLivre}

\clearpage
